% 图契约：并联源支路共用同一端口电压；图中的端口、几何、箭头和方程对应本题。
% 版面：沿用前三模块的几何图+推导分区；目标印宽17cm，简图至中等密度。
% 布局：左侧真实结构，右侧关键关系；真实电路边不作为架构流程箭头处理。
\documentclass[tikz,border=7pt]{standalone}
\usepackage{fontspec,xeCJK,amsmath,amssymb}
\IfFontExistsTF{Times New Roman}{\setmainfont{Times New Roman}}{\setmainfont{TeX Gyre Termes}}
\IfFontExistsTF{Heiti SC}{\setCJKmainfont[AutoFakeBold=1.5]{Heiti SC}}{\setCJKmainfont[AutoFakeBold=1.5]{FandolHei-Regular}}
\usetikzlibrary{arrows.meta,bending,calc,patterns,positioning,decorations.pathmorphing}
\definecolor{ink}{HTML}{243745}
\definecolor{blue}{HTML}{275DAB}
\definecolor{green}{HTML}{177C63}
\definecolor{red}{HTML}{BD4935}
\definecolor{orange}{HTML}{AA6C16}
\definecolor{purple}{HTML}{715393}
\tikzset{every node/.style={font=\fontsize{9.3}{12.5}\selectfont,text=ink},
arrow/.style={-{Stealth[length=5pt 1.5,width'=0pt 3]},shorten >=2pt,shorten <=1pt,line width=1pt},
arrow short/.style={-{Stealth[length=3pt 1.0,width'=0pt 2.5]},shorten >=1pt,shorten <=0pt,line width=.8pt},
axis/.style={arrow,draw=ink!60,line width=.6pt},
vec/.style={arrow,draw=blue},
wire/.style={draw=ink!75,line width=.8pt},
guide/.style={dashed,draw=ink!45,line width=.5pt},
card/.style={rounded corners=2pt,draw=blue!65,fill=blue!4,align=center,text width=6.5cm,inner xsep=9pt,inner ysep=7pt},
title/.style={font=\fontsize{10.5}{14}\selectfont\bfseries,text width=15cm,align=center},
label/.style={font=\fontsize{8.4}{11}\selectfont,align=center},
note/.style={label,text width=6.1cm}}
\begin{document}
\begin{tikzpicture}[x=1cm,y=1cm]
\path[use as bounding box] (-8,-3.5) rectangle(8,3.65);
\node[title] at (0,3.25) {M4-035\quad 并联源支路共用同一端口电压};
\begin{scope}[shift={(-4.05,0)}]
\draw[wire,fill=white](-1.2,1.7)circle(0.3);\node at(-1.2,1.7){};
\node at(-1.34,1.7){$+$};\node at(-1.06,1.7){$-$};\node[above]at(-1.2,2.06){$\mathcal E_1$};\draw[wire](-2.4,1.7)--(-1.5,1.7)(-.9,1.7)--(-.4,1.7);\draw[wire,draw=ink](-0.4,1.7)--(0.725,1.7) (1.275,1.7)--(2.4,1.7);\begin{scope}[shift={(1.0,1.7)},rotate=0.0]\draw[draw=ink,fill=white,line width=.8pt](-0.275,-.085)rectangle(0.275,.085);\end{scope}
\node[label]at(0.9999999999999999,2.12){$r_1$};
\draw[wire,fill=white](-1.2,0)circle(0.3);\node at(-1.2,0){};
\node at(-1.34,0){$+$};\node at(-1.06,0){$-$};\node[above]at(-1.2,0.36){$\mathcal E_2$};\draw[wire](-2.4,0)--(-1.5,0)(-.9,0)--(-.4,0);\draw[wire,draw=ink](-0.4,0)--(0.725,0.0) (1.275,0.0)--(2.4,0);\begin{scope}[shift={(1.0,0.0)},rotate=0.0]\draw[draw=ink,fill=white,line width=.8pt](-0.275,-.085)rectangle(0.275,.085);\end{scope}
\node[label]at(0.9999999999999999,0.42){$r_2$};
\draw[wire,fill=white](-1.2,-1.7)circle(0.3);\node at(-1.2,-1.7){};
\node at(-1.34,-1.7){$+$};\node at(-1.06,-1.7){$-$};\node[above]at(-1.2,-1.3399999999999999){$\mathcal E_3$};\draw[wire](-2.4,-1.7)--(-1.5,-1.7)(-.9,-1.7)--(-.4,-1.7);\draw[wire,draw=ink](-0.4,-1.7)--(0.725,-1.7) (1.275,-1.7)--(2.4,-1.7);\begin{scope}[shift={(1.0,-1.7)},rotate=0.0]\draw[draw=ink,fill=white,line width=.8pt](-0.275,-.085)rectangle(0.275,.085);\end{scope}
\node[label]at(0.9999999999999999,-1.28){$r_3$};
\draw[wire](-2.4,-1.7)--(-2.4,1.7)(2.4,-1.7)--(2.4,1.7);\fill[ink] (-2.4,0)circle(2pt);\node[left=4pt] at(-2.4,0){$A$};
\fill[ink] (2.4,0)circle(2pt);\node[right=4pt] at(2.4,0){$B$};
\node[note]at(0,-2.8){$U=V_A-V_B$；$\mathcal E_i$ 以图示极性取代数值};\end{scope}
\node[card,draw=blue!65,fill=blue!4] at (3.9,1.8) {$\sum_i(U-\mathcal E_i)/r_i=0$};
\node[card,draw=green!65,fill=green!4] at (3.9,0) {$\mathcal E_{\rm eq}=\dfrac{\sum_i\mathcal E_i/r_i}{\sum_i1/r_i}$};
\node[card,draw=orange!65,fill=orange!4] at (3.9,-1.8) {$r_{\rm eq}=(\sum_i1/r_i)^{-1}$};
\end{tikzpicture}\end{document}